Prima di preparare la strumentazione come spiegato nel paragrafo precedente, si è misurata la lunghezza $D$ come somma tra le le lunghezze $ \overline{M_rS_1}+\overline{S_1S_2}+\overline{S_1M_f}$. In particolare, per effettuare tale misura si è utilizzato un metro a nastro dotato di scala graduata con precisione $\sigma_{metro}=1\,cm$. Avendo a che fare con grandezze dell'ordine di metri, ed essendo il nastro morbido, è stato necessario assicurarsi che il nastro rimanesse diritto teso e non si inconcasse al centro. Per fare ciò, uno sperimentatore si è assicurato di tenere sollevata la parte centrale del nastro , evitando allo stesso tempo di modificarne la direzione e cercando di allineare il centro con le estremità.
In questo modo, la misura di $D$ e il suo errore sono state ricavate da
\begin{equation}
    D=\overline{M_rS_1}+\overline{S_1S_2}+\overline{S_1M_f} \, \, \, \, \, \, \, \, \, \, \, 
    \sigma_D=\sqrt{3\sigma_{metro}}=2cm
\end{equation}

Dopo aver sistemato le lenti come spiegato al paragrafo precedente, è stato possibile misurare la grandezza $a$, la distanza tra la seconda lente $L_2$ e lo specchio rotante, grazie al metro presente sul banco ottico, con la precisione di $\sigma_a=1mm$. 

Per le misure di $\Delta\delta$, si è utilizzato il microscopio, dotato di una una croce sulla lente e una rotella sul supporto indispensabili per centrare il punto luminoso corrispondente al raggio di ritorno. 

Una volta scelto il senso di rotazione dello specchio, è stato messo in moto con una frequenza bassa $\nu_0$ (intorno ai 50-100 Hz) e, ruotando appositamente la rotella, si è centrata con la croce del microscopio la posizione in cui appariva il punto luminoso del fascio di ritorno.  
Fatto ciò, si è trovata la misura $\delta_0$ leggendo la scala graduata posta sulla rotella (dotata anche di nonio, in grado di fornire misure con una precisione di 0.01mm)
Dopodiché, si è aumentata la frequenza a un valore $\nu$, si è riposizionato il mirino nel microscopio e si è trovata la misura $\delta_f$. 

Dalla formula (\ref{delta}) si ricava:
\begin{equation}
    c=\frac{4D^2b\omega}{\delta(D+a)}.
    \label{c}
\end{equation},
Sperimentalmente,ci si rende conto che la misura d $b$ risulta estremamente complessa in questo apparato: si sfrutta pertanto la legge dei punti coniugati, per cui si ha
\begin{equation}
    \frac{1}{b}+\frac{1}{D+a}=\frac{1}{f_2}
\end{equation}
dove tutti i valori tranne $b$ sono noti. Se ne ricava
\begin{equation}
    b=\frac{f_2(D+a)}{D+a-f_2}
\end{equation}
che può essere inserito all'interno della formula (\ref{c}) a ottenere:
\begin{equation}
    c=\frac{4f_2D^2(\omega)}{(D+a-f_2)\delta}.
    \label{c'}
\end{equation}
In particolare, risulta significativo partire ogni set di misure con una velocità angolare iniziale non nulla (per ridurre incertezze legate al meccanismo): si sfrutta dunque una simile formulazione delle formula (\ref{c'}) che tiene conto della rotazione iniziale:
\begin{equation}
    c=\frac{4f_2D^2(\omega_f-\omega_0)}{(D+a-f_2)\Delta\delta}
\end{equation}
dove $\omega_f$ e $\omega_0$ rappresentano le velocità angolari relative a $\nu_f$ e $\nu_0$ (si ricorda la legge $\omega=2\pi\nu$).

Una volta effettuate misure per rotazioni dello specchio sia in senso orario che in senso antiorario, è stata utilizzata la funzione di rotazione a massima frequenza e si sono effettuati molteplici misure considerando con $\nu_0$ la massima frequenza in un senso e con $\nu_f$ la massima frequenza nell'altro. In questo modo si è cercato di ottenere un $\Delta\delta$ molto grande.